\documentclass[10pt,a4paper]{amsart}
\usepackage[margin=19mm]{geometry}
\usepackage{amsmath,amssymb,mathtools}
\usepackage[scr=rsfs,cal=euler]{mathalfa}
\usepackage{xfrac}
\usepackage[unicode,hidelinks,bookmarks=false]{hyperref}
\newcommand{\ie}{i.e.\ }
\newcommand{\cf}{cf.\ }
\newcommand{\resp}{respectively\ }
\newcommand{\Subsection}{\subsection}
\providecommand{\nobreakdash}{\nobreak\hbox{-}}
\setlength{\emergencystretch}{2em}
\multlinegap0pt
\usepackage{bm}
\newtheorem{prop}{Proposition}
\title{What aggregation rules can be classified as logical concepts?}
\author{Nikolay L. Poliakov}
\date{Source: arXiv:2604.01260v1 (CC BY 4.0)}
\begin{document}
\maketitle

\noindent\emph{Source and licence.} What aggregation rules can be classified as logical concepts?. Version arXiv:2604.01260v1 (CC BY 4.0), distributed by arXiv under the Creative Commons Attribution 4.0 International licence, http://creativecommons.org/licenses/by/4.0/, as recorded in the arXiv licence metadata for this exact version and shown on the abstract page https://arxiv.org/abs/2604.01260v1. Full licence notice: \texttt{licenses/arxiv-2604.01260-CC-BY-4.0.txt}.\par\smallskip

\noindent\textit{Editorial scope.} Counted unit: the article's prop PropNeutrality (source0.tex lines 192--194). The article's complete original proof of it is delivered with it (source0.tex lines 195--204, one \texttt{proof} environment of 10 lines). No mathematics is written, repaired or re-derived: the statement and the proof are the article's own lines, in the article's own notation, and no retained line is substituted or re-flowed.  No further source-local context is needed: every cross-reference of every retained line resolves inside the two delivered spans.  Nothing was omitted from any retained span, so the package carries no omission record.  No retained line contains a citation, so the wrapper carries no bibliography and no bibliography entry.  No retained line contains a figure environment, an image-inclusion command or drawing code, so no image is delivered and the package declares no asset dependency.  Editorial formatting only, in addition: an AMS \texttt{amsart} preamble that restates the article's own declarations and macros used by the retained lines, namely the environments \texttt{prop} on separate counters, together with the standard packages those lines need.  The retained environments print in this document as Proposition 1 in order of appearance, whereas the article prints the same items with the numbering of its own full document.  The complete native source of the article as distributed in the arXiv e-print for this exact version is bundled unchanged as \texttt{sources/197648/source0.tex}.
\begin{prop}\label{PropNeutrality}
	Every local and neutral aggregation rule $f:C^n\to C$ for $\mathfrak C_2(A)$ preserves any two-element set $D\subseteq C$.
\end{prop}

\begin{proof}
	Let $f:C^n\to C$ be a local and neutral aggregation rule for $\mathfrak C_2(A)$, $D\subseteq C$, $|D|=2$, and $(d^1, d^2, \ldots, d^n)\in D^n$. Choose $c^1, c^2\in D$ and $b\in [A]^2$ such that $c^2(b)\neq c^1(b)=f(d^1, d^2, \ldots, d^n)(b)$. It is enough to show that $f(d^1, d^2, \ldots, d^n)(b')=c^1(b')$ for all $b'\in [A]^2$. If $c^1(b')=c^2(b')$, this follows immediately from the conditions of locality and unanimity. Otherwise, consider a permutation $\sigma\in S_A$ for which $\sigma(c^1(b))=c^1(b')$ and $\sigma(c^2(b))=c^2(b')$. Since $f$ is local, we have:
	$
	f(d^1_\sigma, d^2_\sigma, \ldots, d^n_\sigma)(b)= f(d^1, d^2, \ldots, d^n)(b)=c^1(b).
	$.
	Then since $f$ is neutral, we have:
	$$
	f(d^1, d^2, \ldots, d^n)(b')=\sigma(f(d^1_\sigma, d^2_\sigma, \ldots, d^n_\sigma)(b))=\sigma(c^1(b))=c^1(b').
	$$
\end{proof}

\end{document}
